%% Paste-ready replacement for Section 3 "Testing the Frame" in main.tex.
%% Replaces everything from "Lead time values in our first version" through the
%% paragraph ending "...encoding something real.", including the table.

\section{Testing the Frame}

Lead time values in our first version were conservative manual estimates rather
than statistically computed values. We replaced them with a cross-correlation of
weekly wastewater percentile against weekly clinical cases. Three things came out
of that.

The first is that the Delta wave had to be removed. CDC NWSS percentile reporting
does not begin until 15 November 2021, so the entire Delta window contains no
viral intensity data. Delta had been our flagship wave and it carried most of our
figures. The visualization had been rendering it confidently the whole time.
Nothing in the design forced the metaphor to break, which is the exact failure we
are warning about.

The second is that our strongest numbers came from our weakest data. Ranked on
correlation alone, Idaho led Omicron by three weeks at $r = 0.99$. Idaho has one
sampled treatment plant serving about 2.7\% of the state. We now report a state
only when its sampled plants cover at least 15\% of the population across five or
more sites. Every large lead time we had disappeared under that rule.

The third is that what survives is real but modest (Table~\ref{tab:leads}). Where
a lead exists it is about one week, which sits inside the published range rather
than beyond it. During BA.5 most adequately covered states showed no lead at all,
and West Virginia showed none in either wave.

\begin{table}[t]
  \centering
  \caption{Lead times from cross-correlation of weekly wastewater viral
  percentile against weekly clinical cases, for states where sampled plants cover
  at least 15\% of the population. Coverage is that share. A lead of none means
  the correlation peaks at lag 0.}
  \label{tab:leads}
  \begin{tabular}{llccc}
    \toprule
    State & Wave & Coverage & Lead & $r$ \\
    \midrule
    Illinois      & Omicron & 62\% & 1 week        & 0.90 \\
    Oregon        & Omicron & 51\% & 1 week        & 0.82 \\
    Wisconsin     & Omicron & 23\% & 1 week        & 0.89 \\
    Georgia       & BA.5    & 22\% & 1 week        & 0.92 \\
    Illinois      & BA.5    & 62\% & \textbf{none} & 0.81 \\
    Wisconsin     & BA.5    & 23\% & \textbf{none} & 0.87 \\
    West Virginia & Omicron & 22\% & \textbf{none} & 0.94 \\
    West Virginia & BA.5    & 22\% & \textbf{none} & 0.92 \\
    \bottomrule
  \end{tabular}
\end{table}

This unevenness is not a flaw in the visualization but a reflection of the
real-world variability documented by Kumar et al.~\cite{kumar2022}. We now show
Appalachia and West Virginia rather than leaving them out
(Fig.~\ref{fig:counter}). In those rows the two curves sit on top of each other
and the layers have nothing to separate. The frame fails in public, which is the
strongest evidence we have that it was encoding something real.
